\documentclass{bmstu}
\bibliography{biblio}
\usepackage{graphicx}
\usepackage{subcaption}
\usepackage{amsmath}
\newcommand{\norm}[1]{\lvert #1 \rvert}
\renewcommand{\labelitemi}{---}
\renewcommand{\labelitemii}{$\bullet$}
\renewcommand{\labelitemiii}{$\circ$}

\usepackage{multirow}
\usepackage{booktabs}

\lstset{
	basicstyle=\ttfamily\footnotesize,
	frame=single,
	tabsize=2,
	linewidth=0.95\textwidth,
	xleftmargin=0.025\textwidth,
	numbers=left,
	numbersep=5pt,
	breaklines=true,
	keepspaces=true,
	showspaces=false,
	showstringspaces=false,
	breakatwhitespace=false,
	float=h!,
	abovecaptionskip=-5pt,
}

\begin{document}

\makecourseworktitle
{Информатика и системы управления}
{Программное обеспечение ЭВМ и информационные технологии}
{Разработка компилятора языка Plankalkul}
{Дремин~К.~А./ИУ7-22М}
{Ступников~А.~А.}
{}

\newenvironment{essay2}[1]
{
	\chapter*{РЕФЕРАТ}
	
	Расчетно-пояснительная записка \begin{NoHyper}\pageref{LastPage}\end{NoHyper}~с., \totfig~рис., \tottab~табл., \thetotbibentries~источн., 1~прил.
	
	\MakeUppercase{#1} \par
}{}

\setcounter{page}{3}
\begin{essay2}{Plankalkul, компилятор, ANTLR4, абстрактное синтаксическое дерево, семантический анализ, CIL, CLR, .NET, Python.}
	Работа посвящена разработке компилятора линейной версии языка Plankalkul в промежуточный язык платформы .NET Common Intermediate Language.
	
	В ходе выполнения работы разработано программное обеспечение, выполняющее лексический и синтаксический анализ исходного текста на языке Plankalkul, построение абстрактного синтаксического дерева, семантическую проверку программы и генерацию текстового представления CIL.
	Реализация выполнена на языке \texttt{Python} с использованием генератора синтаксических анализаторов \texttt{ANTLR4}; управление окружением и запуск проверок выполняются через \texttt{uv}.
	
	Разработанная система может применяться для исследования исторических языков программирования, проверки вариантов линейной нотации Plankalkul и демонстрации полного цикла компиляции в исполняемую среду CLR.
\end{essay2}
\maketableofcontents

\chapter*{ВВЕДЕНИЕ}
\addcontentsline{toc}{chapter}{ВВЕДЕНИЕ}

Plankalkul является одним из первых языков программирования высокого уровня.
Язык был предложен Конрадом Цузе для описания алгоритмов независимо от конкретного устройства вычислительной машины.
Несмотря на исторический характер языка, его идеи остаются полезными для изучения ранних подходов к типизации, описанию программ, работе с составными данными и разделению алгоритма на подпрограммы.

Исходная запись Plankalkul была рассчитана на двумерную математическую нотацию.
Для практической реализации компилятора требуется линейное текстовое представление, пригодное для обработки современными средствами синтаксического анализа.
В данной работе рассматривается такая версия языка: она сохраняет основные классы переменных Plankalkul, типы на основе битов, массивы, записи, условные операторы и циклы, но описывается в виде ANTLR4-грамматики.

Целевой платформой компиляции выбран Common Intermediate Language платформы .NET.
Такой выбор позволяет представить каждую программу Plankalkul в виде статического метода CLR, использовать стандартные типы времени выполнения и выполнять сгенерированный код после сборки средствами .NET.

Цель работы -- разработать компилятор языка Plankalkul в CIL.

Для достижения поставленной цели необходимо решить следующие задачи:
\begin{enumerate}
	\item описать грамматику языка Plankalkul, целевую платформу CIL и сформулировать требования к компилятору;
	\item спроектировать компилятор, включающую фронтенд, промежуточное представление, семантический анализ и генерацию CIL;
	\item реализовать обработку программ на языке Plankalkul и интерфейс командной строки;
	\item проверить корректность синтаксического и семантического анализа на тестовых программах.
\end{enumerate}

\chapter{Аналитическая часть}

\section{Язык Plankalkul}
Plankalkul был разработан Конрадом Цузе как язык для записи алгоритмов на уровне, более близком к математическому описанию задачи, чем к машинным командам.
В языке используются переменные разных назначений, явные результаты программ, составные структуры данных и условное выполнение операторов.

В данной работе рассматривается не историческая двумерная запись, а линейная версия языка, основанная на линейной нотации, разработанной и реализованной исследователями Берлинского университета~\cite{berlin-notation}.

\section{Линейная нотация языка}
Линейная нотация представляет программу Plankalkul в виде последовательности текстовых конструкций.
Верхний уровень состоит из двух видов объявлений: определений пользовательских типов и определений программ.
Порядок объявлений свободный: в одном тексте могут последовательно встречаться несколько типов и несколько программ.
Пустая программа также синтаксически допустима, однако для практической компиляции требуется хотя бы одно определение программы.

\subsection{Лексические элементы}
Идентификаторы пользовательских типов и имена программ начинаются со строчной латинской буквы или символа подчеркивания.
После первого символа допускаются латинские буквы в обоих регистрах, цифры и символ подчеркивания.
Числовой литерал записывается десятичными цифрами.
Однобитовое истинное значение задается литералом \texttt{L}, а битовая последовательность-литерал записывается строкой из символов \texttt{L} и \texttt{0}, содержащей хотя бы один символ \texttt{L} и имеющей длину не менее двух символов, например \texttt{L0}, \texttt{LL}, \texttt{00L}.
Комментарии бывают однострочными, начинающимися с \texttt{//}, и блочными, заключенными между \texttt{/*} и \texttt{*/}.

\subsection{Типизация}
Типы в языке задаются типовыми выражениями.
Типовое выражение может использоваться в определении пользовательского типа или как inline-аннотация переменной.
Поддерживаются следующие формы типов:
\begin{itemize}
	\item битовый тип \texttt{0}, описывающий одно логическое значение;
	\item битовая последовательность \texttt{n.0}, где \texttt{n} -- натуральная длина последовательности, например \texttt{8.0};
	\item массив фиксированного размера \texttt{n x T}, где \texttt{n} -- число элементов, а \texttt{T} -- тип элемента;
	\item запись \texttt{(T1, T2, ..., Tn)}, содержащая не менее двух полей;
	\item пользовательский тип, заданный ранее или далее отдельным определением.
\end{itemize}

Определение пользовательского типа имеет форму \texttt{name = type}.
Например, \texttt{int = 8.0} задает восьмиразрядную битовую последовательность, \texttt{node = (int, int)} задает запись из двух полей, а \texttt{node\_pool = 100 x node} задает массив из ста записей.
Типы допускают вложенное построение: элементом массива может быть запись, поле записи может быть массивом, а пользовательский тип может ссылаться на уже построенную составную форму.

Inline-аннотация типа записывается внутри квадратных скобок после двоеточия: \texttt{V0[:int]}, \texttt{Z1[:8.0]}, \texttt{R0[:(int, int)]}.
Такая аннотация используется в сигнатурах программ и при первом явном введении локальных переменных.
В сигнатуре входные параметры и результаты должны иметь явный тип, поскольку по ним формируется тип программы.
Для промежуточных переменных тип может быть указан явно или выведен по правой части присваивания.

\subsection{Структура программы}
Определение программы имеет вид:
\begin{verbatim}
P<number> optional_name(param_list) -> result_spec
  statement_list
END
\end{verbatim}

Идентификатор программы всегда состоит из буквы \texttt{P} и целого номера.
Необязательное имя программы является обычным строчным идентификатором и может использоваться при вызове.

Список параметров заключен в круглые скобки.
Он может быть пустым или состоять из ссылок на переменные, разделенных запятыми.
Параметры задаются переменными класса \texttt{V}; результаты задаются переменными класса \texttt{R}.
Спецификация результата бывает одиночной, например \texttt{R0[:int]}, или множественной, например \texttt{(R0[:int], R1[:int])}.
Тело программы представляет собой последовательность операторов и завершается ключевым словом \texttt{END}.

\subsection{Классы переменных и ссылки на значения}
Переменные делятся на четыре класса:
\begin{itemize}
	\item \texttt{V} -- входные параметры программы;
	\item \texttt{Z} -- промежуточные значения и локальные переменные;
	\item \texttt{C} -- константные значения, допускающие только однократное присваивание;
	\item \texttt{R} -- выходные результаты программы.
\end{itemize}

Ссылка на переменную состоит из класса и номера: \texttt{V0}, \texttt{Z3}, \texttt{C1}, \texttt{R0}.
После номера может следовать квадратная часть, объединяющая компонентный доступ и типовую аннотацию:
\begin{verbatim}
V0[:int]
Z0[1]
Z0[1:node]
V0[(V1 - 1).0:node_pool]
\end{verbatim}

Компонентный доступ состоит из одного или нескольких компонентов, разделенных точками.
Числовой компонент используется для статического выбора поля записи или фиксированного индекса.
Компонент вида \texttt{V1} или \texttt{Z2} задает динамический индекс, взятый из переменной.
Компонент в круглых скобках задает динамический индекс произвольным выражением.
Таким образом, ссылка \texttt{V0[(V1 - 1).0]} означает: взять элемент массива \texttt{V0} с индексом \texttt{V1 - 1}, затем выбрать нулевое поле записи, лежащей в этом элементе.

\subsection{Операторы}
Основной оператор присваивания имеет форму:
\begin{verbatim}
expr => varRef
\end{verbatim}
Левая часть вычисляется как выражение, правая часть должна быть ссылкой на переменную или ее компонент.
Присваивание используется как для записи результата программы, так и для задания промежуточных значений.

Условное выполнение задается guard-оператором:
\begin{verbatim}
expr -> body
\end{verbatim}
Условие должно иметь битовый результат.
Тело guard может быть одиночным оператором или блоком операторов в фигурных скобках.
Так как тело само является оператором, guard допускает вложенные условные конструкции.

Цикл \texttt{W} имеет форму:
\begin{verbatim}
W {
  statement
  ...
}
\end{verbatim}
Его тело должно содержать хотя бы один оператор.
Типичный способ использования \texttt{W} -- размещение внутри одного или нескольких guard-операторов.
Итерации продолжаются, пока на очередном проходе выполняется хотя бы одно условное тело.

Цикл \texttt{FOR} имеет форму:
\begin{verbatim}
FOR varRef FROM expr TO expr {
  statement_list
}
\end{verbatim}
Переменная цикла задается ссылкой на переменную.
Начальное и конечное значения являются выражениями.
Диапазон проходится включительно: тело выполняется для начального значения, затем переменная увеличивается на единицу до достижения конечного значения.

\subsection{Выражения}
Выражения строятся из литералов, ссылок на переменные, вызовов программ, унарных операций, бинарных операций и скобок.
Вызов программы имеет форму \texttt{name(args)} или \texttt{P<number>(args)}.
Список аргументов может быть пустым или состоять из выражений, разделенных запятыми.

Приоритет операций от низкого к высокому следующий:
\begin{itemize}
	\item дизъюнкция \texttt{|};
	\item конъюнкция \texttt{\&};
	\item сравнения \texttt{=}, \texttt{!=}, \texttt{<}, \texttt{>}, \texttt{<=}, \texttt{>=};
	\item сложение и вычитание \texttt{+}, \texttt{-};
	\item умножение и деление \texttt{*}, \texttt{/};
	\item возведение в степень \texttt{**};
	\item унарное отрицание \texttt{\textasciitilde} и унарный минус \texttt{-}.
\end{itemize}

Сравнение в одном выражении не образует цепочку: между двумя арифметическими выражениями допускается не более одного оператора сравнения.
Скобки используются для явного изменения порядка вычисления.
Например, выражение \texttt{V0 + V1 * 2} сначала умножает \texttt{V1} на \texttt{2}, а выражение \texttt{(V0 + V1) * 2} сначала выполняет сложение.

\subsection{Пример программы}
\includelisting
{notation_examples.pk}
{Пример конструкций линейной нотации Plankalkul}

В приведенном примере сначала объявляются пользовательские типы.
Тип \texttt{byte} задает восьмиразрядную битовую последовательность, \texttt{pair} задает запись из двух таких значений, а \texttt{pool} -- массив записей.
Программа \texttt{max} демонстрирует сигнатуру с двумя параметрами и одним результатом, присваивание, guard-оператор и возврат значения через \texttt{R0}.
Программа \texttt{sum} показывает цикл \texttt{FOR}, а \texttt{first\_pair\_field} -- цепочку компонентного доступа к массиву записей.

\section{Целевая платформа}
Целевой платформой компилятора является инфраструктура .NET, а целевым промежуточным представлением -- Common Intermediate Language.
CIL представляет собой стековый язык команд, предназначенный для описания типов, методов, локальных переменных, управляющих переходов и операций над значениями.
После сборки CIL-текст преобразуется в .NET-сборку, которая исполняется средой CLR~\cite{ecma335,dotnet}.

Модель CLR предоставляет набор встроенных примитивных типов, массивы, пользовательские типы-значения, статические методы и точку входа программы.
Эти средства достаточны для представления процедурной линейной нотации Plankalkul: тело программы может быть выражено последовательностью инструкций метода, а операторы управления -- условными и безусловными переходами.
Стековая организация CIL также хорошо соответствует генерации выражений: операнды помещаются на стек, после чего арифметическая, логическая или сравнительная инструкция заменяет их результатом вычисления.

\section{Отображение Plankalkul на целевую платформу}
Программа Plankalkul отображается на статический метод .NET.
Параметры класса \texttt{V} становятся аргументами метода, промежуточные значения \texttt{Z}, константы \texttt{C} и результаты \texttt{R} -- локальными переменными.
Для запуска скомпилированной программы формируется метод \texttt{Main}, который получает аргументы командной строки, преобразует их к целочисленным значениям, вызывает основную программу и выводит результат.

Типы языка отображаются на типы CLR.
Битовые значения представляются типом \texttt{bool}, битовые последовательности длиной до 32 бит -- \texttt{int32}, более длинные последовательности -- \texttt{int64}.
Массивы Plankalkul транслируются в массивы CLR, а записи -- в пользовательские value type со структурными полями.
Такое представление записей сохраняет фиксированную структуру составного значения, а массивы CLR дают возможность работать с индексируемыми областями данных и передавать ссылки на них между методами.

Оператор присваивания отображается на вычисление выражения и последующую запись результата в аргумент, локальную переменную, элемент массива или поле структуры.
Guard-оператор соответствует условному переходу через тело оператора при ложном условии.
Циклы \texttt{W} и \texttt{FOR} выражаются набором меток и переходов: для \texttt{W} проверяется факт срабатывания условных тел на итерации, а для \texttt{FOR} поддерживается переменная цикла и проверка верхней границы диапазона.
Вызов программы Plankalkul отображается на вызов соответствующего статического метода.

\includeimage
	{pipeline}
	{f}
	{H}
	{0.6\textwidth}
	{Общий процесс компиляции Plankalkul в CIL}

\section{Требования к программному обеспечению}
Разрабатываемый компилятор должен:
\begin{itemize}
	\item принимать тексты программ на линейной версии языка Plankalkul;
	\item выполнять лексический и синтаксический анализ для построения абстрактного синтаксического дерева программы;
	\item проверять совместимость типов, корректность условий, циклов и обращений к компонентам;
	\item генерировать текстовый CIL для последующей сборки средствами .NET;
	\item предоставлять команды для разбора, анализа и компиляции из интерфейса командной строки.
\end{itemize}

\section*{Вывод по аналитической части}
В аналитической части рассмотрены особенности линейной версии языка Plankalkul, описана целевая платформа CIL.
Сформулированы основные требования к компилятору, определяющие дальнейшее проектирование системы.

\chapter{Конструкторская часть}

В данном разделе описана архитектура компилятора и основные алгоритмы обработки программы без привязки к деталям конкретной реализации.

\section{Общая архитектура}
Компилятор имеет конвейерную архитектуру.
Исходный текст сначала передается лексеру и парсеру ANTLR4.
Полученное конкретное синтаксическое дерево преобразуется в абстрактное синтаксическое дерево.
Затем выполняется семантический анализ, после которого backend генерирует текст CIL.

\includeimage
	{compiler_components}
	{f}
	{H}
	{\textwidth}
	{Компоненты компилятора Plankalkul}

Основные компоненты системы:
\begin{itemize}
	\item фронтенд, отвечающий за разбор текста и построение AST;
	\item модуль AST, задающий узлы типов, выражений, операторов и программ;
	\item семантический анализатор, формирующий таблицы символов и проверяющий типы;
	\item генератор CIL, преобразующий проверенную программу в текстовую сборку;
	\item управляющий модуль компилятора и интерфейс командной строки.
\end{itemize}

\section{Абстрактное синтаксическое дерево}
AST используется как промежуточное представление программы.
В отличие от конкретного дерева ANTLR, оно не содержит служебных синтаксических узлов и хранит только смысловые элементы языка.

Корневым узлом является программа, содержащая список определений типов и список определений программ.
Определение программы хранит номер \texttt{P}, необязательное имя, список параметров, список результатов и тело.
Тело состоит из операторов присваивания, guarded-операторов, \texttt{W}-циклов и \texttt{FOR}-циклов.
Выражения представлены литералами, ссылками на переменные, вызовами программ, унарными и бинарными операциями.

\section{Семантический анализ}
Семантический анализ выполняется в три фазы.
На первой фазе регистрируются пользовательские типы.
На второй фазе регистрируются сигнатуры всех программ, что позволяет поддерживать вызовы программ, объявленных ниже по тексту.
На третьей фазе проверяются тела программ.

\includeimage
	{semantic_phases}
	{f}
	{H}
	{\textwidth}
	{Фазы семантического анализа}

Для каждой программы формируется область видимости.
Параметры класса \texttt{V} и результаты класса \texttt{R} объявляются по сигнатуре программы.
Промежуточные переменные \texttt{Z} и константы \texttt{C} получают тип при первом использовании с inline-аннотацией или из типа присваиваемого выражения.
Для переменных класса \texttt{C} дополнительно проверяется запрет повторного присваивания.

Типы после разрешения представлены структурно: бит, битовая последовательность, массив и запись.
Совместимость присваивания проверяется по структуре типа.
Для компонентного доступа анализатор последовательно проходит по цепочке компонентов и вычисляет тип выбранного поля записи, элемента массива или бита последовательности.

\section{Генерация CIL}
Каждая программа Plankalkul отображается на статический метод класса \texttt{PlankalkulProgram}.
Параметры класса \texttt{V} становятся аргументами метода.
Переменные \texttt{Z}, \texttt{C} и \texttt{R} размещаются в локальных переменных CIL.
Первый по номеру метод используется как целевая программа для точки входа \texttt{Main}.

Записи отображаются на CLR value type с полями \texttt{f\_0}, \texttt{f\_1} и далее.
Если программа возвращает несколько результатов, генератор создает синтетическую запись и возвращает ее как единое значение.
Массивы создаются инструкцией \texttt{newarr}; доступ к элементам и полям реализуется инструкциями \texttt{ldelem}, \texttt{stelem}, \texttt{ldfld}, \texttt{stfld}.

Арифметические и логические выражения транслируются в стековые инструкции CIL.
Guard-оператор реализуется ветвлением \texttt{brfalse}.
Цикл \texttt{FOR} транслируется в пару меток начала и конца цикла, а \texttt{W}-цикл -- в повторное выполнение тела, пока хотя бы одно условие продолжает выполняться.

\section*{Вывод по конструкторской части}
Была спроектирована модульная архитектура компилятора, разделяющая синтаксический анализ, промежуточное представление, семантическую проверку и генерацию CIL.
Определены основные структуры данных и правила отображения конструкций Plankalkul на элементы CLR.

\chapter{Технологическая часть}

В данном разделе описаны средства разработки и основные особенности реализации компилятора.

\section{Средства разработки}
Компилятор реализован на языке \texttt{Python} версии 3.14.
Управление зависимостями и запуск команд проекта выполняются при помощи инструмента управления проектами \texttt{uv}, что обеспечивает воспроизводимый запуск инструментов из окружения проекта~\cite{python,uv}.

Для синтаксического анализа используется \texttt{ANTLR4} и пакет \texttt{antlr4-python3-runtime}.
Для сборки сгенерированного CIL применяется вспомогательный проект .NET, расположенный в каталоге \texttt{build}.
Проверки проекта выполняются средствами \texttt{pytest}.

\section{Фронтенд компилятора}
Грамматика языка описана в формате \texttt{ANTLR4}.
На ее основе сгенерированы лексер и парсер..
После построения конкретного синтаксического дерева класс построителя AST проходит по узлам ANTLR и создает собственные узлы программы.

Узлы AST реализованы как неизменяемые dataclass-структуры.
Каждый узел содержит позицию в исходном тексте, что позволяет формировать сообщения об ошибках с указанием строки и столбца.
Отдельные группы узлов описывают типы, выражения, операторы, определения типов и определения программ.

\section{Семантический анализ}
Семантический анализатор хранит словари пользовательских типов и программ.
Для каждой программы создается объект области видимости, содержащий известные переменные.
Ошибки накапливаются в списке и выдаются пользователю единым исключением компиляции.

В реализации выделены разрешенные типы: бит, битовая последовательность, массив и запись.
Пользовательский тип при анализе заменяется соответствующим resolved-представлением.
Это позволяет сравнивать типы структурно и не зависеть от имени, которым тип был объявлен в исходной программе.

Особое внимание уделяется компонентному доступу.
При обращении к записи статический компонент выбирает поле по индексу.
При обращении к массиву компонент выбирает тип элемента.
При обращении к битовой последовательности компонент выбирает отдельный бит.
Динамический компонент дополнительно анализируется как выражение.

\section{Генерация CIL}
Генератор CIL формирует текстовую сборку.
Сначала выводится заголовок сборки и определения структур для всех обнаруженных record-типов.
Затем для каждой программы формируется статический метод.
В конце создается метод \texttt{Main}, который разбирает аргументы командной строки как целые числа, вызывает первую программу и выводит результат.

\includelisting
{cil_fragment.il}
{Фрагмент сгенерированного CIL для программы вычисления чисел Фибоначчи}

Для генерации переходов используются уникальные метки.
Такой подход применяется в guard-операторах, \texttt{W}-циклах и \texttt{FOR}-циклах.
Арифметические операции транслируются в инструкции \texttt{add}, \texttt{sub}, \texttt{mul}, \texttt{div}; операция возведения в степень реализуется вызовом \texttt{System.Math::Pow}.

Каждая программа Plankalkul отображается на приватный статический метод CIL.
Если у программы есть имя, оно используется как имя метода; иначе применяется идентификатор вида \texttt{P1}.
Параметры \texttt{V} становятся аргументами метода и загружаются инструкциями \texttt{ldarg}, а промежуточные значения \texttt{Z}, константы \texttt{C} и результаты \texttt{R} помещаются в секцию \texttt{.locals init}.
Для локальных массивов после объявления генерируется явная инициализация: размер помещается на стек инструкцией \texttt{ldc.i4}, затем создается массив инструкцией \texttt{newarr}, после чего ссылка сохраняется в локальную переменную.
Переменные класса \texttt{C} технически также являются локальными переменными CIL, но семантический анализатор запрещает их повторное присваивание.

Отображение типов выполняется после разрешения пользовательских типов.
Бит \texttt{0} соответствует \texttt{bool}, битовые последовательности длиной до 32 разрядов -- \texttt{int32}, более длинные последовательности -- \texttt{int64}.
Массивы языка отображаются на массивы CLR, а записи -- на пользовательские value type структуры.
Для каждой уникальной структурной записи генератор создает тип \texttt{Rec\_N} с полями \texttt{f\_0}, \texttt{f\_1} и далее по числу компонентов.
Именованные типы после семантического анализа сравниваются структурно, поэтому две одинаковые записи могут использовать одно и то же представление в CIL независимо от исходного имени типа.

Результаты программ обрабатываются отдельно от обычных локальных переменных.
Если программа возвращает одно значение, в конце метода загружается соответствующая локальная переменная \texttt{R\_<n>} и выполняется \texttt{ret}.
Если результатов несколько, генератор создает синтетическую record-структуру, заводит локальную переменную \texttt{\_ret}, записывает значения \texttt{R0}, \texttt{R1} и последующих результатов в поля этой структуры инструкциями \texttt{stfld}, затем возвращает структуру целиком.
Метод \texttt{Main} является точкой входа сборки, разбирает аргументы командной строки как \texttt{int32}, вызывает программу с минимальным номером и печатает первый результат через \texttt{System.Console::WriteLine}.

При генерации выражений используется стековая модель CIL.
Целочисленные литералы загружаются инструкциями \texttt{ldc.i4}, булевы значения -- как \texttt{ldc.i4.1} или \texttt{ldc.i4.0}, а битовые последовательности предварительно сворачиваются в целое значение.
Арифметические операции напрямую отображаются на \texttt{add}, \texttt{sub}, \texttt{mul} и \texttt{div}.
Сравнения генерируются через \texttt{ceq}, \texttt{clt}, \texttt{cgt}; операции \texttt{!=}, \texttt{<=} и \texttt{>=} выражаются через строгое сравнение и инверсию результата.
Булевы операции используют \texttt{and} и \texttt{or}, унарное отрицание \texttt{\textasciitilde} реализуется сравнением операнда с нулем, а возведение в степень требует преобразования операндов к \texttt{float64}, вызова \texttt{System.Math::Pow} и обратного преобразования результата к \texttt{int32}.

Компонентный доступ требует различать массивы и записи.
Для массива генератор вычисляет индекс и использует инструкции семейства \texttt{ldelem}/\texttt{stelem}; для массива структур при последующем обращении к полю применяется адресная форма \texttt{ldelema}.
Для записи статический компонент выбирает поле \texttt{f\_<n>}, загрузка выполняется через \texttt{ldfld}, а запись -- через \texttt{stfld}.
При присваивании в компонент генератор сначала получает адрес контейнера, затем вычисляет правую часть и только после этого записывает значение в выбранный элемент или поле.

Управляющие конструкции реализуются через метки и условные переходы.
Guard вычисляет условие и выполняет \texttt{brfalse} на метку конца оператора.
Для цикла \texttt{W} с одним guard используется обычная петля с проверкой условия в начале; для общего случая генератор ведет локальный флаг срабатывания, чтобы повторять цикл только тогда, когда на текущей итерации был выполнен хотя бы один guard.
Цикл \texttt{FOR} инициализирует переменную цикла начальным значением, перед каждой итерацией сравнивает ее с конечным значением, выполняет тело, увеличивает переменную на единицу и переходит к следующей проверке.

\section{Интерфейс командной строки}
Интерфейс командной строки реализует три команды:
\begin{itemize}
	\item \texttt{parse} -- разобрать файл и вывести AST;
	\item \texttt{analyze} -- выполнить синтаксический и семантический анализ;
	\item \texttt{compile} -- сгенерировать CIL и при необходимости собрать программу через .NET.
\end{itemize}

Запуск команд выполняется через \texttt{uv run}, например:
\begin{verbatim}
uv run python -m plankalkulsharp analyze
  tests\programs\test_fibonacci.pk
\end{verbatim}

\section*{Вывод по технологической части}
Была реализована система компиляции на языке \texttt{Python}, использующая ANTLR4 для разбора исходного текста, dataclass-структуры для AST, отдельный семантический анализатор и генератор текстового CIL.
Для воспроизводимого запуска команд проекта используется \texttt{uv}.

\chapter{Исследовательская часть}

В данном разделе описана проверка корректности работы разработанного компилятора на тестовых программах.

\section{Цель исследования}
Цель исследования -- проверить, что компилятор корректно выполняет синтаксический и семантический анализ программ, написанных на текущей версии грамматики Plankalkul, а также генерирует корректный код CIL.

\section{Программа вычисления чисел Фибоначчи}
Первая программа реализует вычисление числа Фибоначчи.
Она демонстрирует определение пользовательского типа, работу с промежуточными переменными, \texttt{W}-цикл, guard-оператор и арифметические выражения.

\includelisting
{test_fibonacci.pk}
{Программа вычисления числа Фибоначчи}

Для проверки использованы команды:
\begin{verbatim}
uv run python -m plankalkulsharp parse
  tests\programs\test_fibonacci.pk
uv run python -m plankalkulsharp analyze
  tests\programs\test_fibonacci.pk
uv run python -m plankalkulsharp compile
  tests\programs\test_fibonacci.pk
  -o build\PlankalkulProgram.il
.\build\bin\Release\net10.0\PlankalkulProgram.exe 10
\end{verbatim}

Синтаксический анализ завершился построением AST.
Семантический анализ завершился без ошибок и предупреждений.
В отчете анализатора зарегистрирован тип \texttt{int} как битовая последовательность длиной 8 и программа \texttt{P1 fib} с одним параметром и одним результатом.
После команды \texttt{compile} был создан файл \texttt{build\textbackslash PlankalkulProgram.il} и собран исполняемый файл \texttt{build\textbackslash bin\textbackslash Release\textbackslash net10.0\textbackslash PlankalkulProgram.exe}.
Запуск полученной программы с аргументом \texttt{10} вывел значение \texttt{55}, что соответствует десятому числу Фибоначчи при нумерации, используемой в тестовой программе.

\section{Программа разворота связного списка}
Вторая программа демонстрирует более сложный набор возможностей языка.
В ней используются пользовательские типы \texttt{int}, \texttt{pointer}, \texttt{node}, \texttt{node\_pool}, массив записей, несколько программ, вызовы программ, \texttt{FOR}-цикл, \texttt{W}-цикл и компонентный доступ.

\includelisting
{test_reverse_list_fragment.pk}
{Фрагмент программы разворота связного списка}

Для проверки использованы команды:
\begin{verbatim}
uv run python -m plankalkulsharp parse
  tests\programs\test_reverse_list.pk
uv run python -m plankalkulsharp analyze
  tests\programs\test_reverse_list.pk
uv run python -m plankalkulsharp compile
  tests\programs\test_reverse_list.pk
  -o build\PlankalkulProgram.il
.\build\bin\Release\net10.0\PlankalkulProgram.exe 5 5
\end{verbatim}

Синтаксический анализ завершился построением AST.
Семантический анализ завершился без ошибок и предупреждений.
В отчете анализатора зарегистрированы записи, массив записей и программы \texttt{P0 test}, \texttt{P1 alloc\_pool}, \texttt{P4 create\_list}, \texttt{P5 reverse\_list}.
После компиляции была собрана та же исполняемая точка входа \texttt{PlankalkulProgram.exe}, но уже с CIL-кодом программы разворота списка.
Запуск с аргументами \texttt{5 5} вывел \texttt{4}: для списка длины 5 после разворота узел со значением 5 указывает на узел со значением 4, что проверяет изменение связей между элементами массива записей.

\section{Результаты проверки}
В таблице \ref{tbl:research_results} приведены результаты проверки актуальных программ.

\includelisting
{exe_test_output.txt}
{Фрагмент консольного вывода при компиляции и запуске актуальных программ}

\begin{table}[H]
	\caption{Результаты проверки актуальных программ}
	\label{tbl:research_results}
	\begin{tabular}{|p{4.8cm}|p{3.8cm}|p{4.7cm}|}
		\hline
		Программа & Проверяемые возможности & Результат \\ \hline
		\texttt{test\_fibonacci.pk} & Типы, присваивания, \texttt{W}-цикл, guard, арифметика & Разбор и анализ без ошибок; сборка exe успешна; запуск с аргументом \texttt{10} выводит \texttt{55} \\ \hline
		\begin{tabular}[t]{@{}l@{}}\texttt{test\_reverse}\\\texttt{\_list.pk}\end{tabular} & Записи, массивы, компонентный доступ, вызовы программ, \texttt{FOR}, \texttt{W} & Разбор и анализ без ошибок; сборка exe успешна; запуск с аргументами \texttt{5 5} выводит \texttt{4} \\ \hline
	\end{tabular}
\end{table}


\section*{Вывод по исследовательской части}
В результате исследования показано, что тестовые программы на языке Plankalkul успешно проходят синтаксический и семантический анализ, компилируются в CIL и собираются в исполняемый файл .NET.
Запуск полученных программ подтвердил работоспособность сгенерированного кода для арифметического цикла и для операций над массивом записей, моделирующим связный список.

\chapter*{ЗАКЛЮЧЕНИЕ}
\addcontentsline{toc}{chapter}{ЗАКЛЮЧЕНИЕ}
В ходе выполнения работы была достигнута поставленная цель -- разработано программное обеспечение, реализующее метод кодогенерации инструкций доступа к данным в графовом хранилище данных из SQL-выражений.

В процессе выполнения были решены все поставленные задачи:
\begin{enumerate}
	\item Проведен анализ существующих подходов к кодогенерации и особенностей предметной области;
	\item Спроектирована библиотека кодогенерации инструкций для работы с графовой базой данных;
	\item Разработана библиотека кодогенерации и интерфейс командной строки к ней;
	\item Исследована корректность генерации инструкций работы с графовой базой данных из SQL-выражений.
\end{enumerate}

В процессе исследования было выяснено, что разработанное решение успешно позволяет генерировать инструкции доступа к данным в составе вычислительного комплекса <<Тераграф>> на языке \texttt{C++} из SQL-выражений.

\phantomsection
\addcontentsline{toc}{chapter}{СПИСОК ИСПОЛЬЗОВАННЫХ ИСТОЧНИКОВ}
\printbibliography[title=СПИСОК ИСПОЛЬЗОВАННЫХ ИСТОЧНИКОВ]

\begin{appendices}
\iffalse
	\chapter{}
	Приведены листинги примитивных объектов и некоторых комбинаторов.
	\includelistingpretty
	{Map/Objects/Impl/Box.cs}
	{[Sharp]C}
	{Класс параллелепипеда}
	\includelistingpretty
	{Map/Objects/Impl/Cylinder.cs}
	{[Sharp]C}
	{Класс сферы}
    \includelistingpretty
    {Map/Objects/Impl/Plane.cs}
    {[Sharp]C}
    {Класс плоскости}
    \includelistingpretty
    {Map/Objects/Impl/Sphere.cs}
    {[Sharp]C}
    {Класс сферы}
    \includelistingpretty
    {Map/Objects/Combinators/ObjectUnion.cs}
    {[Sharp]C}
    {Класс объединения}
	\includelistingpretty
	{Map/Objects/Combinators/ObjectDifference.cs}
	{[Sharp]C}
	{Класс разности}
	\includelistingpretty
	{Map/Objects/Combinators/ObjectIntersection.cs}
	{[Sharp]C}
	{Класс пересечения}
	\includelistingpretty
	{Map/Objects/Combinators/ObjectRoundUnion.cs}
	{[Sharp]C}
	{Класс скруглённого объединения}
	\chapter{}
	Приведены листинги реализаций некоторых моделей освещения.
	\includelistingpretty
	{Map/Light/LambertLight.cs}
	{[Sharp]C}
	{Класс модели освещения по Ламберту}
	\includelistingpretty
	{Map/Light/PhongLight.cs}
	{[Sharp]C}
	{Класс модели освещения Фонга}
	\includelistingpretty
	{Map/Light/BlinnLight.cs}
	{[Sharp]C}
	{Класс модели освещения Блинна-Фонга}
	\includelistingpretty
	{Map/Light/BlinnSoftAoLight.cs}
	{[Sharp]C}
	{Класс модели освещения Блинна-Фонга с мягкими тенями и окружающим затенением}
	\chapter{}
	Приведены листинги кодогенераторов основных функций отрисовки кроме генератора модели освещения, рассмотренного в технологической части.
	\includelistingpretty
	{Map/Generators/MapGenerator.cs}
	{[Sharp]C}
	{Класс кодогенератора ФРЗ}
	\includelistingpretty
	{Map/Generators/MaterialGenerator.cs}
	{[Sharp]C}
	{Класс кодогенератора функции определения материала}
	\includelistingpretty
	{Map/Generators/TransformGenerator.cs}
	{[Sharp]C}
	{Класс кодогенератора преобразований координат}
	\includelistingpretty
	{Map/Generators/UniformGenerator.cs}
	{[Sharp]C}
	{Класс кодогенератора изменяемых извне переменных шейдера}
	\includelistingpretty
	{Map/Generators/FragmentGenerator.cs}
	{[Sharp]C}
	{Класс кодогенератора итогового фрагментного шейдера}
	\chapter{}
	Приведены листинги привязки программы на GLSL к программе на \texttt{C\#}.
	\includelistingpretty
	{Render/RenderScene.cs}
	{[Sharp]C}
	{Класс настройки отображения сцены на экране}
	\includelistingpretty
	{Render/Shader.cs}
	{[Sharp]C}
	{Класс взаимодействия с шейдером}
\fi
	\chapter{}
	Приведена презентация к курсовой работе
\end{appendices}

\end{document}
